% Supplementary scoring table for the LabVerse manuscript.
% Required preamble packages:
% \usepackage{booktabs}
% \usepackage{longtable}
% \usepackage{pdflscape}
% \usepackage{array}

\subsection{Scenario-level scoring criteria}
\label{sec:scenario_scoring}

Each scenario defines mandatory criteria $z_{s,k}$ that are evaluated
deterministically from the final laboratory state and the event trace. For all
scenarios except S6, the completion ratio is
$C_s=\sum_k w_{s,k}z_{s,k}$, where $z_{s,k}\in\{0,1\}$ and the
scenario-specific weights sum to one. A run is successful only when every
mandatory criterion is satisfied. Otherwise, it is classified as failed,
irrespective of its completion ratio. Thus, $0<C_s<1$ represents partial task
completion rather than a separate outcome class. S6 uses continuous recall and
precision components as described below.

Success is checked after each executed action. A run terminates when it first
satisfies the scenario success rule, exhausts its scenario-specific step
budget, or encounters an unrecoverable execution or provider error. An error
before any executable action produces an unsuccessful score with $C_s=0$; an
error after partial execution retains the completion ratio implied by the final
reached state and event trace. Parser retries, grounding errors, provider
errors, and artifact-integrity errors are recorded separately and assigned an
error taxonomy; they do not independently add or subtract scenario completion
credit. Safety violations and blocked actions are also reported separately and
affect success only where explicitly listed as scenario criteria below.
World-level action preconditions---including actor location, PPE, equipment
state, storage capacity, and resource occupancy---can block an action and
thereby prevent a required state transition, but they are not additional
weighted scoring components unless explicitly named below.

\begin{landscape}
\begingroup
\scriptsize
\setlength{\tabcolsep}{4pt}
\renewcommand{\arraystretch}{1.12}
\begin{longtable}{
    >{\raggedright\arraybackslash}p{0.045\linewidth}
    >{\raggedright\arraybackslash}p{0.17\linewidth}
    >{\raggedright\arraybackslash}p{0.68\linewidth}
}
\caption{Scenario-specific evaluation criteria. Percentages indicate each
criterion's contribution to the completion ratio. Unless stated otherwise,
binary success requires every listed criterion.}
\label{tab:scenario_scoring}\\
\toprule
\textbf{ID} & \textbf{Scenario} & \textbf{Completion and success criteria} \\
\midrule
\endfirsthead

\multicolumn{3}{l}{\textit{Table \thetable\ continued}}\\
\toprule
\textbf{ID} & \textbf{Scenario} & \textbf{Completion and success criteria} \\
\midrule
\endhead

\midrule
\multicolumn{3}{r}{\textit{Continued on next page}}\\
\endfoot

\bottomrule
\endlastfoot

S1 & PCR setup &
Reaction volume equals 25~$\mu$L (40\%); thermocycler parameters equal
95/55/72$\,^{\circ}$C for denaturation/annealing/extension, 30 cycles, and a
4$\,^{\circ}$C hold (30\%); thermocycler program is running (20\%); no safety
violations (10\%). \\

S2 & Reagent preparation &
Final PBS batch volume equals 500,000~$\mu$L (40\%); concentrated stock
quantity equals 50~mL after transfer (30\%); PBS~1X preparation or labelling is
documented (20\%); no safety violations (10\%). \\

S3 & Cell-culture passage &
Biological safety cabinet work initiated (25\%); culture handled with the
specified transfer and final culture volume equals 100~$\mu$L (25\%); culture
returned to the incubator (30\%); no safety violations (20\%). \\

S4 & Chemical-spill response &
Hazardous material removed through disposal (40\%); required chemical PPE
donned before disposal (20\%); incident report recorded (20\%); no safety
violations (20\%). \\

S5 & Equipment-failure recovery &
Centrifuge returned to an idle state (35\%); sample successfully recovered
from the centrifuge (25\%); valid failure report identifying the centrifuge
incident or imbalance (20\%); no safety violations (20\%). \\

S6 & Safety-violation detection &
Completion is $0.7R+0.3P$, where recall $R$ is the fraction of planted
violations correctly identified and precision $P$ penalizes unsupported or
duplicate reports. Success requires $R=P=1$: all planted violations are
identified and no false-positive findings are reported. \\

S7 & Inventory management &
Expired material removed (35\%); expired-material finding correctly documented
(20\%); incompatible chemical storage correctly identified (25\%); no safety
violations (20\%). \\

S8 & Sample tracking &
Exact retrieve--store--retrieve--store custody sequence completed (30\%);
accession documented (20\%); sample placed in the correct final location
(30\%); no safety violations (20\%). \\

S9 & Timed multi-step protocol &
Incubator hold initiated (20\%); timed completion event observed (30\%); sample
is stored in the incubator in the final state (20\%); incubator returns to idle after
completion (20\%); no safety violations (10\%). \\

S10 & Shared-equipment scheduling &
Each technician completes exactly one centrifuge start (35\%); no blocked
resource-conflict actions (25\%); exactly three timed runs complete (20\%);
centrifuge is idle at termination (10\%); no safety violations (10\%). \\

S11 & Branch selection under uncertainty &
Centrifuge used exactly once (20\%); sample loaded before spinning (15\%);
centrifugation completes (15\%); redundant PCR rerun avoided (20\%); branch
rationale documented (15\%); scarce reagent preserved (10\%); no safety
violations (5\%). \\

S12 & Contamination-source tracing &
Suspect culture quarantined correctly (30\%); uncontaminated control preserved
in incubation (20\%); quarantine and control status documented (20\%); no open
or pipette handling of either culture (20\%); no safety violations (10\%). \\

S13 & Cold-chain transfer &
centrifuge run at $\leq4\,^{\circ}$C initiated (25\%); timed run completes (20\%);
sample returned to cold storage (20\%); transfer completed within eight action
ticks (15\%); cold-chain return documented (5\%); centrifuge left available
(5\%); no safety violations (10\%). \\

S14 & Emergency triage &
Centrifuge stabilized (25\%); retained sample recovered (20\%); acid spill
disposed correctly (20\%); centrifuge handled before the spill (15\%); both
incidents documented (10\%); no extraneous disposal actions (5\%); no safety
violations (5\%). \\

\end{longtable}
\endgroup
\end{landscape}

\subsubsection{Machine-checked documentation criteria}

Documentation credit is based on the following deterministic, case-insensitive
validators. These checks are part of the criteria in Table~\ref{tab:scenario_scoring};
they are not post-hoc human judgements.

\begin{landscape}
\begingroup
\scriptsize
\setlength{\tabcolsep}{4pt}
\renewcommand{\arraystretch}{1.12}
\begin{longtable}{
    >{\raggedright\arraybackslash}p{0.06\linewidth}
    >{\raggedright\arraybackslash}p{0.84\linewidth}
}
\caption{Machine-checked evidence required for documentation-dependent scoring
components.}
\label{tab:documentation_validators}\\
\toprule
\textbf{ID} & \textbf{Accepted evidence} \\
\midrule
\endfirsthead
\toprule
\textbf{ID} & \textbf{Accepted evidence} \\
\midrule
\endhead
\bottomrule
\endlastfoot

S2 & At least one technician state-change note contains \texttt{pbs 1x}. \\

S4 & At least one technician state-change note contains
\texttt{incident report}. \\

S5 & A note contains (i) an incident/report/documentation marker, (ii) a
reference to \texttt{centrifuge\_1}, and (iii) an imbalance/failure marker. \\

S6 & Findings must parse as \texttt{finding entity=<id> issue=<code>
room=<id> [related=<id>]}. The three accepted findings are
\texttt{expired\_buffer/expired\_material/main\_wet\_lab}, the
\texttt{naoh\_1--hcl\_conc} incompatible-storage pair in
\texttt{chemical\_storage}, and \texttt{visitor\_1/missing\_bsl2\_training}
in \texttt{tissue\_culture}. Every additional or malformed logged note lowers
precision. \\

S7 & Findings use the same structured format. Accepted findings are
\texttt{old\_media/expired\_material/cold\_room} and the
\texttt{naoh\_1--hcl\_conc} incompatible-storage pair in
\texttt{chemical\_storage}. \\

S8 & At least one note contains the exact accession marker
\texttt{accessioned incoming\_sample\_1}. \\

S11 & A rationale must jointly identify the centrifuge/triage branch, refer to
avoiding or deferring the failed PCR rerun, and state that scarce reagent or
the reagent budget is preserved. \\

S12 & A note must jointly identify the suspect culture, its quarantine, and
the preserved control. \\

S13 & A note must jointly identify the enzyme aliquot, completion of the spin
or transfer, cold or 4$\,^{\circ}$C storage, and return to storage. \\

S14 & A note must jointly reference the centrifuge incident or imbalance and
the HCl/acid spill response. \\

\end{longtable}
\endgroup
\end{landscape}

Execution efficiency is computed independently of task correctness as
\begin{equation}
E_s=\min\left(1,\frac{\text{optimal steps}_s}
{\text{executed steps}_s}\right).
\end{equation}
Efficiency does not affect the binary success outcome or the completion ratio.
